\section{Related Work}
\label{sec:related_work}

% ---------------------------------------------------------------------------
\subsection{Garment Pattern Nesting}
\label{sec:rw_nesting}

The problem of placing irregularly shaped panels on a strip to minimise
wasted material is known as the irregular strip packing or nesting
problem~\cite{WascherHS07}.
Bennell and Oliveira~\cite{BennellO08} provide a comprehensive tutorial
covering the principal geometric tools, from raster methods to the
no-fit polygon (NFP), and Le{\~a}o et al.~\cite{LeaoTOCA20} give a recent
review of mathematical models for irregular packing.
The NFP, which characterises all valid non-overlapping placements of one
polygon relative to another, underpins many state-of-the-art nesting
solvers~\cite{BurkeHKW07}.
Burke et al.~\cite{BurkeHKW06} introduce the bottom-left-fill (BLF)
heuristic, which places each piece at the lowest, leftmost feasible position
on the strip; Mart{\'i}nez-Sykora et al.~\cite{MartinezSykoraACRO22}
recently proposed a fast, scalable BLF variant based on a semi-discrete
representation.
Our nesting engine builds on the classical BLF paradigm.

Metaheuristic approaches have been widely applied to nesting.
Jakobs~\cite{Jakobs96} introduces one of the first genetic algorithm
formulations for polygon packing, encoding placement order as a permutation.
Mundim et al.~\cite{MundimAQO17} apply a biased random-key genetic
algorithm to open-dimension nesting with a no-fit raster representation.
Heckmann and Lengauer~\cite{HeckmannL95} address the textile industry
specifically with simulated annealing, and
Gomes and Oliveira~\cite{GomesO06} combine simulated annealing with linear
programming for local compaction.
Elkeran~\cite{Elkeran13} achieves strong results with guided cuckoo search
and pairwise clustering.

All of the above methods treat texture as irrelevant: panels are placed
purely for geometric compactness with no regard for pattern continuity
across seams.
NFP-based solvers represent the state of the art for fabric efficiency
alone; our nesting component uses a BLF heuristic since the novelty lies
in the joint objectives, and any NFP-based engine could be substituted
without affecting the other contributions.

% ---------------------------------------------------------------------------
\subsection{Texture and Pattern Alignment}
\label{sec:rw_texture}

Ko and Kim~\cite{Ko2013} address nesting on figured (patterned) fabric using
image analysis and 3D drape simulation to preserve figure continuity across
pieces.
Their approach jointly considers placement and pattern matching, but is
limited to simple repeat patterns and does not handle the full symmetry
structure of periodic textures.

Wolff and Sorkine-Hornung~\cite{WolffS19} present the most comprehensive
prior work on seam-aware texture alignment in garments.
They formalise the alignment problem using all 17 wallpaper groups and
optimise the placement of sewing pattern pieces to minimise phase mismatch
along seams, with optional deformation of the seam shapes themselves.
A follow-up~\cite{WolffHS19} extends this framework to enforce body
reflection symmetry constraints on both texture alignment and seam shape.
While this line of work establishes the formal foundations for wallpaper
group alignment, it assumes an unbounded fabric supply and ignores material
efficiency entirely.
Furthermore, by operating directly on 2D patch geometry, any deformation
introduced to improve alignment implicitly modifies the garment fit, with
no mechanism to measure or bound this change.

Pietroni et al.~\cite{PietroniDFLVS22} address a complementary problem:
given a 3D garment shape, they segment it into panels and compute 2D
parameterizations that respect tailoring constraints such as seam symmetry,
dart placement, and fabric grain alignment, with a flattening distortion
measure that models woven-fabric deformation.
Their objective is to optimise seam placement for minimal flattening
distortion, whereas our method optimises seam placement for texture
alignment under fabric roll constraints.


% ---------------------------------------------------------------------------
\subsection{Multi-Objective Evolutionary Algorithms and Mixed-Type Genotypes}
\label{sec:rw_moea}

The seminal NSGA-II algorithm~\cite{Deb2002NSGAII} establishes the standard
framework for Pareto-based multi-objective optimisation with tournament
selection: candidate solutions are compared by dominance, with a
crowding-distance tiebreaker to maintain diversity.
Deb and Jain~\cite{DebJ14} extend this to many-objective problems in
NSGA-III with reference-point-based selection.
We adopt a simplified variant of the NSGA-II selection scheme, using
fitness sum as the tiebreaker in place of crowding distance, which is
sufficient for our three-objective problem.

The mixed structure of our genotype (continuous seam parameters, discrete
orientations and phase offsets, and a combinatorial nesting order) requires
dedicated operators per variable type.
Michalewicz~\cite{Michalewicz1996} establishes the foundational treatment
of mixed-type chromosomes in evolutionary computation, arguing that each
variable type benefits from operators that respect its domain structure.
Li et al.~\cite{LiEEBSDR13} formalise this principle in Mixed Integer
Evolution Strategies (MIES), defining separate self-adaptive mutation
operators for continuous, integer, and nominal discrete variables within
a unified algorithm.
Our design follows this principle, as detailed in Section~\ref{sec:ea}.
